%% kapa: exam=1 round=29 year=2018 period=2 form=A subject=accounting subject_ko=회계학 seq=27 number=67
%% source: Q-Net 감정평가사 기출문제 (cst00309) artlSeq=5202808 file=extracted/1008/5202808/2018년 제29회 감정평가사 1차시험 기출문제/2018년 제29회 감정평가사 1차시험 기출문제/2교시 A형.pdf
%% conversion: kordoc 4.24.51 -> tools/gen1.py (auto); review=proofread-sonnet; flags=none
\begin{question}{67}
재고자산에 관한 설명으로 옳지 않은 것은?
\begin{choices}
  \item 재료원가, 노무원가 및 기타 제조원가 중 비정상적으로 낭비된 부분은 재고자산의 취득원가에 포함시키지 않고 발생기간의 비용으로 인식한다.
  \item 제작기간이 단기간인 재고자산은 차입원가를 자본화 할 수 있는 적격자산에 해당되지 아니한다.
  \item 매입거래처로부터 매입수량이나 매입금액의 일정률만큼 리베이트를 수령할 경우 이를 수익으로 인식하지 않고 재고자산 매입원가에서 차감한다.
  \item 기말 재고자산은 취득원가와 순실현가능가치 중 낮은 금액으로 측정한다.
  \item 가격변동이익이나 중개이익을 목적으로 옥수수, 구리, 석유 등의 상품을 취득하여 단기간 내에 매도하는 기업은 순공정가치의 변동을 기타포괄손익으로 인식한다.
\end{choices}
\end{question}
\answer{5}
